
\appendix
\chapter{Richiami di Analisi e Serie}

Questa appendice contiene nozioni fondamentali di analisi matematica necessarie per lo studio dei modelli probabilistici, in particolare limiti superiore/inferiore e serie.

\section{Liminf e Limsup di successioni reali}

Sia $(x_n)_{n \in \N}$ una successione di numeri reali. 
La successione dell'estremo inferiore definitivo è definita come:
\[ y_n := \inf_{k \ge n} x_k \]
Questa successione $(y_n)$ è sempre crescente. Di conseguenza, ammette sempre limite (finito o infinito).

\begin{definition}[Limite Inferiore]
Il \textbf{limite inferiore} di una successione $(x_n)$ è definito come:
\[ \liminf_{n \to \infty} x_n := \lim_{n \to \infty} \left( \inf_{k \ge n} x_k \right) \]
\end{definition}

In maniera simmetrica, definendo la successione degli estremi superiori definitivi $z_n := \sup_{k \ge n} x_k$, che è sempre decrescente:

\begin{definition}[Limite Superiore]
Il \textbf{limite superiore} di una successione $(x_n)$ è definito come:
\[ \limsup_{n \to \infty} x_n := \lim_{n \to \infty} \left( \sup_{k \ge n} x_k \right) \]
\end{definition}

\begin{spiegazione}[A cosa servono $\liminf$ e $\limsup$?]
Il limite normale $\lim x_n$ potrebbe non esistere (pensa a $x_n = (-1)^n$, che oscilla tra -1 e 1 in eterno). Tuttavia, $\liminf$ e $\limsup$ esistono \textbf{sempre}, perché misurano i confini della successione a lungo termine. 
Se il limite normale esiste, allora $\liminf$ e $\limsup$ sono uguali, e coincidono col limite!
\end{spiegazione}

In generale, vale la disuguaglianza:
\[ \liminf_{n \to \infty} x_n \le \limsup_{n \to \infty} x_n \]
E si ha l'uguaglianza $\liminf x_n = \limsup x_n$ se e solo se il limite $\lim x_n$ esiste.

\section{Serie}

\subsection{Serie a termini non negativi}

Sia $I$ un insieme di indici (es. numerabile, con la stessa cardinalità di $\N$) e $f : I \to [0, +\infty]$. Data una biiezione $\sigma : \N \to I$, definiamo la somma parziale:
\[ S_n := \sum_{k=0}^n f(\sigma(k)) \]
Poiché $f \ge 0$, la successione $(S_n)$ è crescente, e perciò ammette limite (finito o infinito). Questo limite:
\[ \lim_{n \to \infty} \sum_{k=0}^n f(\sigma(k)) = \sum_{x \in I} f(x) \]
\textbf{non dipende} dalla scelta della biiezione $\sigma$. Ovvero, in una serie a termini positivi, puoi riordinare gli addendi come preferisci, il risultato finale sarà sempre lo stesso.

\subsection{Serie assolutamente convergenti}

Cosa succede se i termini possono essere negativi? Sia $f : I \to \R$.
Supponiamo che la serie dei valori assoluti converga:
\[ \sum_{x \in I} |f(x)| < \infty \]
In tal caso, la serie si dice \textbf{assolutamente convergente}. Anche in questo caso, è possibile dimostrare che il limite della somma:
\[ \sum_{x \in I} f(x) \]
esiste finito, e soprattutto \textbf{non dipende dall'ordine in cui si sommano i termini}.

\begin{hint}[Attenzione all'ordine nelle serie!]
Il Teorema di Riemann sulle serie dice che se una serie converge, ma non assolutamente (es. $\sum (-1)^n / n$), riordinando opportunamente i termini puoi farla convergere a \emph{qualsiasi} numero reale tu voglia! Per questo l'assoluta convergenza (o il segno costante) è cruciale per poter maneggiare le serie come se fossero somme normali.
\end{hint}

\chapter*{Bibliografia}
\addcontentsline{toc}{chapter}{Bibliografia}
\begin{itemize}
    \item [1] D. Acemoglu, \emph{Introduction to Modern Economic Growth}, Princeton University Press, 2009.
    \item [2] Q. Berger, F. Caravenna, and P. Dai Pra, \emph{Probabilità}, Springer, 2021.
    \item [3] D. Bertsekas, \emph{Dynamic Programming and Optimal Control}, Athena Scientifica, 2005.
    \item [5] P. Billingsley, \emph{Probability and Measure}, Wiley, 1995.
    \item [6] T. Bjork, \emph{Arbitrage theory in continuous time}, Oxford University Press, 2009.
    \item [7] P. Brémaud, \emph{Markov Chains}, Springer, 1999.
    \item [13] A. Pascucci and W. Runggaldier, \emph{Finanza Matematica: Teoria e Problemi per Modelli Multi-periodali}, Springer, 2009.
\end{itemize}
